\subsection{四元数代数的中心}

设 \(\alpha, \beta, \gamma\) 是 \(K\) 的元素，\(F\) 是型为 \((\alpha, \beta, \gamma)\) 的四元数代数。

命题 2. —— a) 设域 K 的特征异于 2。若 \(\gamma\) 或 \(4\alpha + \beta^{2}\) 非零，则 F 的中心等于 K；否则，它的维数为 2，且由 1 与 ij - ji 生成。

\begin{enumerate}
\def\labelenumi{\alph{enumi})}
\setcounter{enumi}{1}
\tightlist
\item
  设域 \(\mathrm{K}\) 的特征为 2。若 \(\beta \neq 0\)，则 \(\mathrm{F}\) 的中心等于 \(\mathrm{K}\)；若 \(\beta = 0\)，则代数 \(\mathrm{F}\) 是交换的。
\end{enumerate}

由 III, §2, No.~5, 第 444 页的公式 (30)，我们有

\[
i j - j i = - \beta j + 2 k, \quad j k - k j = \beta \gamma - 2 \gamma i, \quad k i - i k = - 2 \alpha j - \beta k. \tag {4}
\]

F 的元素 \(q = x + yi + zj + tk\) 是中心的，当且仅当它与 i 和 j 交换，即我们有

\[
2 z + \beta t = - \beta z + 2 \alpha t = 0 \quad \text { and } \quad 2 \gamma t = \beta \gamma t = 2 y = \beta y = 0.\tag{5}
\]

先设 K 的特征异于 2。若 \(\gamma\) 非零，则 (5) 中的等式蕴含 y = t = 0，进而 z = 0；于是我们有 \(q \in K\)。若 \(\gamma = 0\) 且 \(4\alpha + \beta^{2} \neq 0\)，则它们蕴含 y = z = t = 0，于是我们有 \(q \in K\)。若 \(\gamma = 4\alpha + \beta^{2} = 0\)，则方程组 (5) 约化为 \(y = 2z + \beta t = 0\)，于是我们有 \(q = x + t/2(ij - ji)\)，这就完成了 a) 的证明。

现设域 K 的特征为 2。这时方程组 (5) 可写成 \(\beta t = \beta z = \beta y = 0\)；断言 b) 随之得出。

